\paragraph{4. Comets.}
In this section we define comets and prove a comet-finding lemma which will be used in our sampling lower bound.
\paragraph{Definition 20.}
A $d$-comet is a triple of indices $(i,j,k)$ from $[m]$ with $i<j<k$ such that $j-i\geq d(k-j)$. We call $(j,k)$ the head and $(i,j)$ the tail. The head length is $k-j$. A set of comets $\{(i_h,j_h,k_h)\}_h$ is disjoint if the intervals $[i_h,k_h]$ are disjoint.
\paragraph{Lemma 21 (Comet-finding).}
A subset of $\{1,2,\ldots,m\}$ of size $m/\ell^b$ contains $\geq m/\ell^{b+c+O(1)}$ disjoint $\ell^c$-comets where the head lengths are all in $[\ell^h,\ell^{h+1}]$ for some integer $h\leq b+c+O(1)$, for any $m$, $b\leq\ell$, $c\leq\ell$, and $\ell\geq\log m$.
\begin{proof}
Let $d=\ell^c$. First we claim that any subset of size $n:=d\log m+2$ contains a $d$-comet. Let the elements in the set be $a_1,a_2,\ldots$ in increasing order. If $(a_1,a_2,a_3)$ is not a $d$-comet then $a_3-a_2>(a_2-a_1)/d$, and so $a_3-a_1=a_3-a_2+a_2-a_1\geq(a_2-a_1)(1+1/d)$. Then again if $(a_1,a_3,a_4)$ is not a $d$-comet we have $a_4-a_3\geq(a_3-a_1)/d$ and so $a_4-a_1\geq(a_3-a_1)(1+1/d)\geq(a_2-a_1)(1+1/d)^2$. If we continue this way $n-2$ times, we obtain $a_n\geq(1+1/d)^{n-2}>m$, which is a contradiction.
Now divide the $t:=m/\ell^b$ elements of the given set into consecutive blocks of size $n$. By the previous paragraph, each block contains a comet. Hence we have $\geq t/n-1$ disjoint $d$-comets.
At least half of these comets have heads of length $\leq O(mn/t)=\ell^{b+c+O(1)}$, otherwise half the comets have heads longer than that, and we run out of space. Let $C_i$ be the subset of these comets whose head length is in $[\ell^i,\ell^{i+1})$. We only need to consider $i\leq b+c+O(1)$. Hence, there exists $i=h$ and
\[\Omega\left(\frac tn\right)\frac1{b+c+O(1)}\geq\frac m{\ell^{b+c+O(1)}}\]
disjoint comets with head lengths in $[\ell^h,\ell^{h+1})$, using that both $b$ and $c$ are $\leq\ell$.
\end{proof}
\paragraph{5. A lemma about entropy.}
In this section we quickly recall a basic result about entropy which will be used in our sampling lower bounds. The entropy $H$ of a random variable $X$ is defined as $H(X):=\sum_x\Pr[X=x]\cdot\lg(1/\Pr[X=x])$. The conditional entropy $H(X\mid Y):=\mathbb E_{y\in Y}H(X\mid Y=y)$ (cf. Chapter 2 in [13]).
\paragraph{Lemma 22.}
Let $Z=(Z_1,\ldots,Z_k)$ where $Z_i$ is supported over a set $S_i$, and let $\sum_i\log|S_i|=M$. Suppose $H(Z)\geq M-a$. There is a set $G\subseteq[k]$ of size $|G|\geq k-a/\epsilon$ such that for any $i\in G$ we have
\[H(Z_i\mid Z_1Z_2\ldots Z_{i-1})\geq\log|S_i|-\epsilon.\]
In particular, $Z_i$ is $4\sqrt\epsilon$ close to uniform over $S_i$.
\begin{proof}
By the chain rule for entropy ([13], Equation 2.21)
\[\sum_{i\leq k}(\log|S_i|-H(Z_i\mid Z_1Z_2\ldots Z_{i-1}))\leq a.\]
Applying Markov inequality to the non-negative random variable $\log|S_i|-H(Z_i\mid Z_1Z_2\ldots Z_{i-1})$ (for random $i\in[k]$), we have
\[\mathbb P_{i\in[k]}[\log|S_i|-H(Z_i\mid Z_1Z_2\ldots Z_{i-1})\geq\epsilon]\leq a/(k\cdot\epsilon),\]
yielding the desired $G$.
The ``in particular'' part holds because conditioning reduces entropy: $H(Z_i)\geq H(Z_i\mid Z_1Z_2\ldots Z_{i-1})$ ([13], Equations 2.60 and 2.92) and then applying Pinsker's inequality ([14], Chapter 3; Exercise 17).
\end{proof}
